\subsection{Retrieval-Augmented Generation}
Retrieval-Augmented Generation (RAG) has been widely applied to enhance Large Language Models (LLMs) by retrieving relevant information from external sources, addressing restricted context windows, improving factuality, and mitigating hallucinations~\cite{fan2024survey, gao2023retrieval}. 
Most RAG systems process text corpora by splitting documents into chunks~\cite{finardi2024chronicles}. 
Given a query, relevant chunks can be retrieved via lexical search~\cite{ram2023context} or semantic similarity search~\cite{karpukhin2020dense}. 
Beyond vanilla retrieval, pre-retrieval processing~\cite{ma2023query, zheng2023take}, post-retrieval processing~\cite{dong2024don, xu2023recomp}, and fine-tuning strategies~\cite{li2023structure} further improve effectiveness across tasks such as question answering~\cite{yan2024corrective}, dialogue generation~\cite{izacard2023atlas}, and summarization~\cite{jiang2023active}. 
Many systems also employ reranking and iterative retrieval to refine evidence selection and improve answer quality under a fixed context budget. Several studies benchmark RAG pipelines and evaluation tools across tasks and domains~\cite{yu2024evaluation, chen2024benchmarking, es2023ragas}. 
However, existing work rarely provides a controlled comparison between standard RAG and GraphRAG under unified experimental settings on widely used text benchmarks.

\vspace{-0.1in}
\subsection{Graph Retrieval-Augmented Generation}
Many real-world scenarios involve graph-structured data, such as knowledge graphs (KGs), social graphs, and molecular graphs~\cite{xia2021graph, ma2021deep}. 
GraphRAG incorporates graph structures into retrieval to exploit relational signals among connected nodes~\cite{han2024retrieval, peng2024graph}. 
Early work primarily studies retrieval over existing KGs for downstream tasks such as KG-based QA~\cite{tian2024graph, yasunaga2021qa} and fact checking~\cite{kim2023factkg}. 
Graph structures can also benefit text-centric retrieval; for example, hyperlink graphs between documents can improve retrieval for question answering~\cite{li2022dynamic}.
Recent works further explore constructing graphs from text to support text-based tasks~\cite{han2024retrieval}. 
One direction builds document- or chunk-level graphs to guide retrieval over textual units~\cite{dong2024don, li2022dynamic}. 
Another direction constructs entity--relation graphs from documents (often with LLM assistance) and retrieves information at multiple abstraction levels, such as local neighborhoods or community-level summaries~\cite{edge2024local, han2025reasoning}. 
More recently, text-centric and hierarchical approaches use graph-inspired structures without fully specified entity--relation semantics: RAPTOR constructs hierarchical summary structures for multi-granular retrieval~\cite{sarthi2024raptor}, while HippoRAG and its extensions build entity-linked graphs to guide chunk retrieval~\cite{jimenez2024hipporag, gutierrez2025rag}. 

Despite rapid progress, GraphRAG systems for text are often evaluated under heterogeneous protocols, with varying graph construction methods, retrieval configurations, and even evaluation criteria. 
Moreover, graph construction introduces additional costs (indexing time, retrieval latency, and storage footprint) and can be sensitive to the quality of the construction model, yet these trade-offs are not consistently characterized across studies. 
As a result, it remains unclear how GraphRAG compares with standard RAG on general text-based benchmarks and what practical trade-offs it entails. This motivates our systematic benchmark evaluation under unified experimental settings.
